\chapter[Degree-Five Technical Ledger]{Technical Ledger for the Degree-Five Calculation}

This appendix collects the exact structure used in Chapters 5 and 6 so that the long-form argument can remain readable without sacrificing the algebraic certificate. The representation is the real irreducible $SO(3)$ module $V_5^{\mathbb R}$ of dimension eleven, restricted to the unit sphere $S^{10}$. The baseline parameters used in the quartic reduced-energy calculation are $\mu=1/5$, $\xi=1/50$, and $\nu=g=1$.

\section*{Invariant multiplicities through degree six}

An independent exact Molien computation reproduces

\[
\dim\operatorname{Inv}_d(V_5)=1,0,1,0,2,0,6
\qquad(d=0,\ldots,6),
\]

with the raw symmetric-power dimensions satisfying

\[
\dim\operatorname{Sym}^d(V_5)=\binom{d+10}{10}.
\]

Thus

\[
\dim\operatorname{Inv}_2=1,
\qquad
\dim\operatorname{Inv}_4=2,
\qquad
\dim\operatorname{Inv}_6=6.
\]

The zero multiplicities at degrees three and five are specific to $V_5$ and must not be read as a universal odd-degree rule for $SO(3)$ invariants.

If $I_2(A)=\|A\|^2$, multiplication by $I_2$ is injective, so

\[
\dim(I_2\operatorname{Inv}_4)=2,
\]

and therefore

\[
\boxed{
\dim\left(\operatorname{Inv}_6/I_2\operatorname{Inv}_4\right)=4.
}
\]

\section*{Quartic channel relations and reduced energy}

The quartic channel identities are

\[
K_4=\frac{18}{13}K_0,
\]

\[
K_6=\frac{45}{17}K_0-\frac{143}{170}K_2,
\]

\[
K_8=\frac{880}{247}K_0-\frac{117}{95}K_2,
\]

\[
K_{10}=\frac{777}{323}K_0+\frac{693}{646}K_2.
\]

Hence

\[
F_4(A)=\alpha K_0(A)+\beta K_2(A),
\]

with

\[
\alpha=\frac{30241998291659}{24945292961264\pi},
\qquad
\beta=\frac{503273130591}{3050785612808\pi}>0.
\]

On $S^{10}$, $K_0=1/11$, and with $B_2(A)=P_2(A\otimes A)$,

\[
F_4\big|_{S^{10}}=\frac\alpha{11}+\beta\|B_2(A)\|^2.
\]

Therefore

\[
\mathcal M_4=\{A\in S^{10}:B_2(A)=0\}.
\]

The explicit point

\[
A_\star=\sqrt{\frac35}Y_{50}+\sqrt{\frac25}Y^c_{55}
\]

satisfies $B_2(A_\star)=0$. The differential is

\[
DB_2(A)[\delta A]=2P_2(A\otimes\delta A).
\]

An explicit five-by-five minor of $DB_2(A_\star)$ has determinant

\[
\frac{224\sqrt{1430}}{2924207}\neq0,
\]

so the differential has rank five. Since the target $V_2$ is five-dimensional, this is maximal rank. The radial direction belongs to the kernel because $DB_2(A_\star)[A_\star]=2B_2(A_\star)=0$, hence the tangent restriction retains rank five and the implicit-function theorem yields

\[
\dim T_{A_\star}\mathcal M_4=5.
\]

Let $R_x,R_y,R_z$ denote the infinitesimal $SO(3)$ generators acting on $A_\star$. Their Gram matrix is

\[
R^\top R=10I_3,
\]

so the orbit tangent has dimension three. Equivariance of $B_2$ implies $DB_2(A_\star)R_i=0$. The residual tangent space

\[
N_{\mathrm{shape}}=\ker(DB_2|_T)\cap T_{A_\star}(SO(3)\cdot A_\star)^\perp
\]

therefore has dimension two. One convenient orthogonal spanning pair is

\[
U=-\frac{\sqrt{14}}{14}Y^c_{52}+Y^c_{53},
\]

\[
V=\frac{\sqrt{14}}{14}Y^s_{52}+Y^s_{53},
\]

with

\[
\|U\|^2=\|V\|^2=\frac{15}{14},
\qquad
\langle U,V\rangle=0.
\]

\section*{Explicit sextic witness}

Define

\[
B_4(A)=P_4(A\otimes A),
\qquad
C_8(A)=P_8(B_4(A)\otimes A),
\]

and

\[
J_{4,8}(A)=\|C_8(A)\|^2.
\]

An independent construction from the same spherical-harmonic and Clebsch--Gordan conventions reproduces

\[
J_{4,8}(A_\star)=\frac1{143},
\]

\[
DJ_{4,8}(A_\star)[U]=DJ_{4,8}(A_\star)[V]=0.
\]

\section*{Intrinsic versus ambient second derivatives}

For intrinsic second derivatives along the constrained quartic minimum manifold, a path through $A_\star$ in tangent direction $U$ is written

\[
A(\epsilon)=A_\star+\epsilon U+\frac12\epsilon^2W_U+O(\epsilon^3).
\]

The unit-sphere condition gives

\[
\langle A_\star,W_U\rangle=-\|U\|^2,
\]

while the second-order $B_2=0$ condition gives

\[
DB_2(A_\star)[W_U]+2P_2(U\otimes U)=0.
\]

One exact pair of corrections is

\begin{align*}
W_U={}&-\frac{7\sqrt{15}}{30}Y_{50}
+\frac{\sqrt{21}}{21}Y^c_{51}
+\frac{\sqrt{21}}7Y^c_{54}
-\frac{13\sqrt{10}}{70}Y^c_{55},\\
W_V={}&-\frac{7\sqrt{15}}{30}Y_{50}
-\frac{\sqrt{21}}{21}Y^c_{51}
-\frac{\sqrt{21}}7Y^c_{54}
-\frac{13\sqrt{10}}{70}Y^c_{55}.
\end{align*}

A valid mixed correction is

\[
W_{UV}=\frac{\sqrt{21}}{21}Y^s_{51}-\frac{\sqrt{21}}7Y^s_{54}.
\]

The ambient second derivative is

\[
D^2J_{4,8}(A_\star)[U,U]=\frac{547}{1001},
\]

so the normalized straight-line tangent contribution is

\[
\frac{14}{15}\frac{547}{1001}=\frac{1094}{2145}.
\]

This is not the constrained Hessian. The intrinsic second derivative is

\[
\frac{d^2}{d\epsilon^2}J(A(\epsilon))\Big|_0
=D^2J(A_\star)[U,U]+DJ(A_\star)[W_U],
\]

with

\[
DJ_{4,8}(A_\star)[W_U]=-\frac{43}{1001}.
\]

Hence

\[
\frac{547}{1001}-\frac{43}{1001}=\frac{72}{143}.
\]

The exact intrinsic derivatives are therefore

\[
d^2_{\mathcal M_4}J_{4,8}(U,U)
=d^2_{\mathcal M_4}J_{4,8}(V,V)
=\frac{72}{143},
\]

\[
d^2_{\mathcal M_4}J_{4,8}(U,V)=0.
\]

With

\[
e_1=\sqrt{\frac{14}{15}}U,
\qquad
e_2=\sqrt{\frac{14}{15}}V,
\]

one obtains

\[
\boxed{
H_{\mathrm{shape}}(J_{4,8})=\frac{336}{715}I_2\succ0.
}
\]

The difference between $1094/2145$ and $336/715$ is therefore geometric: the first is the normalized ambient-Hessian contribution, while the second includes the second-fundamental-form correction required to stay on $S^{10}\cap\{B_2=0\}$.

\section*{Frozen theorem and open dynamical step}

The degree-six algebra contains four primitive directions, and the explicit sextic invariant $J_{4,8}$ resolves both nonrotational quartic-flat directions with positive intrinsic curvature. This is an exact algebraic resolving-capability theorem.

The missing dynamical step is the fifth-order amplitude or equivalent sixth-order reduced-energy calculation for the original nonlinear PDE. That reduction must determine which primitive sextic invariants actually appear and with what coefficients. In particular, it must decide the restricted matrix

\[
H_{6,\mathrm{shape}}^{\mathrm{PDE}}
=
\begin{pmatrix}
 d^2H_6(e_1,e_1) & d^2H_6(e_1,e_2)\\
 d^2H_6(e_2,e_1) & d^2H_6(e_2,e_2)
\end{pmatrix}_{\mathcal M_4}.
\]

No claim in this book treats the existence or verified curvature of $J_{4,8}$ as a substitute for that PDE calculation.
